% Using exesheet as a PACKAGE on top of plain article, not as its own
% document class: exesheet.cls additionally \RequirePackage{schooldocs}
% for page layout/headers, which is not needed here and is not in this
% site's pinned TeX Live -- the .sty form has no such dependency.
\documentclass[a4paper]{article}
% hmarginratio/marginparwidth: exesheet puts point values and grading
% notes in the margin via \marginpar, which needs a wide enough margin to
% actually render (the default article/geometry margin is too narrow) --
% following the class's own documented \geometry{} recipe.
\usepackage[a4paper,hmarginratio=2:1,marginparwidth=2.5cm]{geometry}
\usepackage[english]{babel}
\usepackage{amsmath}
% output=questions typesets a clean student copy; switch to
% output=answers for a correction/grading copy with the worked answers
% shown instead. display=notes shows both the point value AND the
% grading-comment text (from \note{}) in the margin; display=pts shows
% only the point value.
\usepackage[output=questions,display=notes]{exesheet}

\begin{document}

\begin{center}
  {\Large \textbf{Chemistry --- Worksheet 5: Stoichiometry}}\\[0.3em]
  Show all working. Molar masses: H = 1, C = 12, O = 16 g/mol.
\end{center}

\bigskip

\exercise[\points{10}]
\question{Balance the following combustion reaction:}
\[
  \mathrm{C_3H_8 + O_2 \longrightarrow CO_2 + H_2O}
\]
\answer{$\mathrm{C_3H_8 + 5\,O_2 \longrightarrow 3\,CO_2 + 4\,H_2O}$.
Balancing carbon gives 3 CO$_2$; balancing hydrogen gives 4 H$_2$O; that
fixes 10 oxygen atoms on the right, so 5 O$_2$ on the left.}

\subpart[\points{4}]
\question{How many moles of O$_2$ are needed to fully combust
$2\,\text{mol}$ of C$_3$H$_8$?}
\answer{From the balanced equation, the mole ratio of C$_3$H$_8$ to
O$_2$ is $1:5$, so $2\,\text{mol}$ of C$_3$H$_8$ needs
$10\,\text{mol}$ of O$_2$.}

\subpart[\points{4}]
\question{What mass of CO$_2$ is produced from that same
$2\,\text{mol}$ of C$_3$H$_8$?}
\answer{The mole ratio of C$_3$H$_8$ to CO$_2$ is $1:3$, so
$2\,\text{mol}$ of C$_3$H$_8$ produces $6\,\text{mol}$ of CO$_2$.
The molar mass of CO$_2$ is $44\,\text{g/mol}$, so the mass produced is
$6 \times 44 = 264\,\text{g}$.}

\note{Common error: forgetting to balance oxygen last, after carbon and
hydrogen are already fixed.}

\exercise[\points{15}]
\question{A solution is made by dissolving $5.85\,\text{g}$ of NaCl
(molar mass $58.5\,\text{g/mol}$) in enough water to make
$500\,\text{mL}$ of solution.}

\subpart[\points{5}]
\question{Calculate the number of moles of NaCl dissolved.}
\answer{$n = m/M = 5.85 / 58.5 = 0.1\,\text{mol}$.}

\subpart[\points{5}]
\question{Calculate the molarity of the solution.}
\answer{$c = n/V = 0.1\,\text{mol} / 0.5\,\text{L} = 0.2\,\text{mol/L}$.}

\subpart[\points{5}]
\question{If the solution is diluted to $2\,\text{L}$, what is the new
molarity?}
\answer{Moles of solute are unchanged on dilution:
$c_2 = n/V_2 = 0.1\,\text{mol} / 2\,\text{L} = 0.05\,\text{mol/L}$.}

\annex
This worksheet assumes familiarity with balancing chemical equations and
basic mole-ratio stoichiometry from Chapter 4.

\end{document}
